2. \underline{数值结果}

2.1. 设 $\underline a=(a_i)_{1\leq i\leq d}$ 为由 $\geq1$ 的整数构成的序列, $V_n(\underline a)$ 表示域 $k$ 上射影空间 $\mathbb P^{n+d}$ 中维数为 $n$、多重次数为 $\underline a$ 的光滑完全交. 令

\begin{equation*}
h^{p,q}(\underline a)
=\dim_k H^q\bigl(V_{p+q}(\underline a),\Omega^p\bigr)
\end{equation*}

以及

\begin{equation*}
h_0^{p,q}(\underline a)=h^{p,q}(\underline a)-\delta_{p,q}.
\end{equation*}

于是 $h_0^{p,q}(\underline a)$ 是多重次数为 $\underline a$ 的完全交的中间维本原上同调的 Hodge 数. 由 1.5, 它们只依赖于 $\underline a$, $p$ 和 $q$. 置

\begin{equation}\tag{2.1.1}
H(\underline a)=\sum_{p,q\geq0}h_0^{p,q}(\underline a)y^pz^q
\in\mathbb Z[[y,z]].
\end{equation}

当 $d=1$ 且 $\underline a=(a)$ 时, 置 $h^{p,q}(a)=h^{p,q}(\underline a)$ 和 $H(a)=H(\underline a)$. 当 $d=0$ 时, 置 $\sup(a_i)=1$.

以下公式是 Hirzebruch [3] 第~160 页公式 (2) 在 $k=0$ 时的特例:

%%%%%%%%%% source scan idx 59 | folio 52 %%%%%%%%%%

\begin{equation}\tag{2.2}
\sum_{n=0}^{\infty}\sum_p
\chi\bigl(V_n(\underline a),\Omega^p\bigr)y^pz^{n+d}
=\frac{1}{(1+yz)(1-z)}
\prod_{i=1}^d
\frac{(1+yz)^{a_i}-(1-z)^{a_i}}
{(1+yz)^{a_i}z+y(1-z)^{a_i}}.
\end{equation}

我们将其变换为如下公式.

\underline{定理} 2.3 (Hirzebruch). \underline{沿用 2.1 的记号, 有}

\begin{equation*}
H(\underline a)=\frac{1}{(1+y)(1+z)}\left[
\prod_{i=1}^d
\frac{(1+y)^{a_i}-(1+z)^{a_i}}
{-(1+y)^{a_i}z+(1+z)^{a_i}y}-1\right].
\end{equation*}

由 1.5,

\begin{equation*}
\chi\bigl(V_n(\underline a),\Omega^p\bigr)
=(-1)^p+(-1)^{n-p}h_0^{p,n-p}(\underline a).
\end{equation*}

因此, (2.2) 的左端 $I$ 可改写为

\begin{align*}
I
&=\sum_{n=0}^{\infty}\sum_{p+q=n}
\left((-1)^p+(-1)^q h_0^{p,q}(\underline a)\right)y^pz^{p+q+d}\\
&=\sum_{p,q\geq0}(-1)^p(yz)^pz^qz^d
+\sum_{p,q\geq0}(-1)^q h_0^{p,q}(\underline a)(yz)^pz^qz^d\\
&=z^d\left[\frac{1}{(1+yz)(1-z)}
+H(\underline a)(yz,-z)\right].
\end{align*}

(2.2) 的右端 $II$ 可改写为

\begin{equation*}
II=\frac{z^d}{(1+yz)(1-z)}
\prod_{i=1}^d
\frac{(1+yz)^{a_i}-(1-z)^{a_i}}
{(1+yz)^{a_i}z+yz(1-z)^{a_i}}.
\end{equation*}

在 $I=II$ 中约去 $z^d$, 得到

\begin{equation*}
H(\underline a)(yz,-z)
=\frac{1}{(1+yz)(1-z)}\left[
\prod_{i=1}^d
\frac{(1+yz)^{a_i}-(1-z)^{a_i}}
{(1+yz)^{a_i}z+(1-z)^{a_i}yz}-1\right],
\end{equation*}

作变量代换即得 2.3.

%%%%%%%%%% source scan idx 60 | folio 53 %%%%%%%%%%

\underline{推论} 2.4. \underline{沿用 2.1 的记号:}

\begin{enumerate}[label=(\roman*),leftmargin=2.6em]
\item \underline{当 $d=1$ 时, 有}
\begin{equation*}
\begin{aligned}
H(a)
&=-\frac{(1+y)^{a-1}-(1+z)^{a-1}}
{(1+y)^az-(1+z)^ay}\\[1ex]
&=\frac{\displaystyle\sum_{i,j\geq0}
\binom{a-1}{i+j+1}y^iz^j}
{\displaystyle1-\sum_{i,j\geq1}\binom{a}{i+j}y^iz^j}.
\end{aligned}
\end{equation*}

\item \underline{以 $|P|$ 表示 $[1,d]$ 的子集 $P$ 的元素个数, 则有}
\begin{equation*}
H(\underline a)=
\sum_{\substack{P\subset[1,d]\\P\ne\varnothing}}
(1+y)^{|P|-1}(1+z)^{|P|-1}
\prod_{i\in P}H(a_i).
\end{equation*}
\end{enumerate}

由 2.3,

\begin{align*}
H(a)
&=-\frac{1}{(1+y)(1+z)}\left[
\frac{(1+y)^a-(1+z)^a}{(1+y)^az-(1+z)^ay}+1\right]\\
&=-\frac{1}{(1+y)(1+z)}\left[
\frac{(1+y)^a(1+z)-(1+z)^a(1+y)}
{(1+y)^az-(1+z)^ay}\right]\\
&=-\frac{(1+y)^{a-1}-(1+z)^{a-1}}
{(1+y)^az-(1+z)^ay}.
\end{align*}

设 $N$ 和 $D$ 分别为最后一个表达式中的分子与分母, 因而 $H(a)=-N/D$.

%%%%%%%%%% source scan idx 61 | folio 54 %%%%%%%%%%

有

\begin{equation*}
N=\sum_k\binom{a-1}{k}(y^k-z^k)
=(y-z)\sum_{i,j\geq0}\binom{a-1}{i+j+1}y^iz^j
\end{equation*}

以及

\begin{align*}
D
&=\sum_k\binom ak(y^kz-yz^k)
=(z-y)+yz\sum_{k\geq1}\binom ak(y^{k-1}-z^{k-1})\\
&=(y-z)\left(-1+yz\sum_{i,j\geq0}
\binom a{i+j+2}y^iz^j\right)\\
&=(y-z)\left(-1+\sum_{i,j\geq1}
\binom a{i+j}y^iz^j\right).
\end{align*}

这就证明了 (i).

现在将 2.3 在 $d=1$ 时的公式反解, 得到

\begin{equation*}
-\frac{(1+y)^a-(1+z)^a}{(1+y)^az-(1+z)^ay}
=1+(1+y)(1+z)H(a).
\end{equation*}

将此恒等式代入 2.3, 得到

\begin{align*}
H(\underline a)
&=\frac{1}{(1+y)(1+z)}\left[
\prod_{i=1}^d\left(1+(1+y)(1+z)H(a_i)\right)-1\right]\\
&=\frac{1}{(1+y)(1+z)}\left[
\sum_{\substack{P\subset[1,d]\\P\ne\varnothing}}
(1+y)^{|P|}(1+z)^{|P|}\prod_{i\in P}H(a_i)\right],
\end{align*}

这就证明了 (ii).
